% ================================
% Chapter 1 — Introduction & Motivation (2–page target)
% Paste after you switch to \pagenumbering{arabic}
% ================================
% ==========================================================
\chapter{Introduction \& Motivation}

\section{Introduction}
The Advanced Encryption Standard (AES) remains the most widely adopted block cipher in
modern secure systems, supporting key lengths of 128, 192, and 256 bits. Its suitability for
hardware stems from a regular round-based structure, enabling efficient implementation on
FPGA and ASIC platforms. A variety of baseline AES architectures have been developed on
reconfigurable hardware, demonstrating functional correctness and area–throughput tradeoffs
\cite{sai_srinivas_fpga_2016, abhijith_high_2013}. Work on compact lookup-table based
transformations further reduced redundant hardware in S-Box and inverse S-Box circuits
\cite{rajeswara_rao_fpga_2017}. Systematic FPGA benchmarking environments such as
\cite{guruprasad2018evaluation} have contributed to consistent evaluation of AES hardware
cores across platforms.

Despite the algorithm's structural simplicity, conventional AES designs incur considerable
area and dynamic power due to duplicated round logic, continuous toggling of SubBytes and
MixColumns units, and the overhead of key schedule expansion. Optimized architectures have
attempted to address these issues through resource sharing and datapath restructuring.
Notably, the Virtex-7 optimized AES-256 design reported in 
\cite{gunasekaran_aes256_virtex7} introduced a unified S-Box/Inv-S-Box structure, a
word-serial datapath, and a combined MixColumns--AddRoundKey unit, achieving significant
reductions in hardware footprint. Such improvements align with early efforts to minimize
resource usage in high-performance FPGA realizations of AES
\cite{abhijith_high_2013, rajeswara_rao_fpga_2017}.

In recent years, power efficiency has emerged as a key design objective due to the growing
presence of AES engines in mobile, IoT, and battery-constrained environments. Studies such
as \cite{feng2021_energy_efficient_aes, ibrahim2021_highspeed_resourceefficient_aes}
demonstrate that even optimized datapaths can remain power-intensive when active every
cycle. As a result, selective activation techniques—primarily clock gating—have gained
traction. Fine-grained asynchronous activation methods proposed in
\cite{selvapriya2023lowpower} exhibit substantial savings by enabling only the logic
required in a given cycle. Similarly, gated-datapath approaches in 
\cite{darwish2022configurable} target MixColumns and SubBytes units to reduce unnecessary 
switching. Multi-domain gating strategies reported in \cite{lee2021multidomain} further 
demonstrate that isolating clock domains for S-Box, MixColumns, and KeySchedule blocks can 
achieve meaningful power reduction.

Direct application of clock gating to AES functional blocks has also been validated through 
FPGA experiments. In \cite{chodorowski2021_clockgated_aes}, simple gating of non-active
cycles led to measurable savings in switching power, while pipelined AES designs using
clock gating on 28 nm technology achieved notable improvements in energy per operation
\cite{prasanth2021_lowpower_aes_clockgating}. Additional S-Box-centric gating mechanisms 
have been explored in \cite{subramanian2021_adaptive_clockgated_sbox}, highlighting the 
benefits of turning off composite-field logic when unused.

Motivated by these developments, this thesis constructs a fully functional baseline AES
implementation supporting all three key sizes. The primary contribution is a systematic,
dual-architecture study: a \textbf{full-pipeline AES implementation} and a \textbf{finite
state machine (FSM) AES implementation} are both developed, characterized at baseline, and
subjected to two optimization phases. Phase~1 replaces GF multiplication using EXP3 and LN3
log-table lookups with \textit{xtime} arithmetic~\cite{wolkerstorfer_xtime_2002}, eliminating
512 table entries from every MixColumns instance across both architectures. Phase~2 introduces
clock enable guards on FSM state registers to suppress unnecessary switching when the datapath
is idle. The Virtex-7 AES-256 architecture of \cite{gunasekaran_aes256_virtex7} provides
structural context for the unified S-Box and combined MixColumns--AddRoundKey design choices
adopted in the FSM datapath. The pipeline implementation follows the parallel datapath
approach established in prior work~\cite{good_pipelined_aes_2005}, enabling high throughput
through unrolled round logic.


\section{Motivation}

Modern VLSI and FPGA systems place strict constraints on power consumption, silicon area, and resource utilization. Cryptographic accelerators such as AES are now integrated into a wide range of platforms including communication modules, embedded processors, cloud servers, and edge devices. These systems require secure and high throughput encryption while also maintaining low energy consumption and minimal hardware cost. Conventional AES architectures, although functionally correct, often rely on duplicated S Box hardware, full width datapaths, and continuously active functional units. This leads to considerable dynamic power usage and limits the practicality of deployment in energy constrained systems.

The increasing use of AES 256 further amplifies these challenges due to its larger number of rounds and the additional complexity of the KeySchedule unit. This results in higher switching activity, increased latency, and greater pressure on available FPGA resources. Therefore, there is a strong need for an AES architecture that reduces redundant computation, improves hardware reuse, and manages switching activity intelligently.

Recent research demonstrates that clock gating is one of the most effective techniques for reducing dynamic power in cryptographic hardware. Fine grained gating, as explored in \cite{selvapriya2023lowpower}, selectively activates functional units only when required. Gated datapath techniques reported in \cite{darwish2022configurable} show that disabling MixColumns and SubBytes during idle cycles can produce substantial power savings. Multi domain gating methods in \cite{lee2021multidomain} further validate the potential for controlling switching activity across independent AES components. These works highlight a clear opportunity to improve AES efficiency by combining structural optimizations with targeted power aware design.

In addition, the evaluation framework proposed in \cite{guruprasad2018evaluation} emphasizes the importance of systematic performance analysis for cryptographic hardware. While this framework does not focus on power minimization, it underscores the need for reproducible and structured comparison of architectural choices. Motivated by these insights, this work aims to design an AES architecture that is both area efficient and energy aware. By integrating hardware reuse techniques inspired by \cite{gunasekaran_aes256_virtex7} with selective clock gating strategies from recent low power research, the goal is to develop a scalable AES core that is optimized for modern FPGA and VLSI platforms.

A central goal is to quantify the area and power tradeoff between the two architectural
paradigms: the pipeline architecture~\cite{good_pipelined_aes_2005}, which provides high
throughput through full round unrolling, and the FSM
architecture~\cite{mcloone_fsm_aes_2003}, which achieves area efficiency through round
hardware reuse. The two optimization phases applied equally to both architectures provide a
controlled, reproducible comparison of these design tradeoffs.

